\chapter{Sampling, Quantization and Digital Telephony}
\label{ch:pcm}

\begin{chapterintro}
Every digital communication system begins where the analog world ends: at a sampler and a
quantizer. This chapter explains what is lost and what is kept when a continuous waveform
becomes a sequence of numbers. It develops the sampling theorem and its practical cousins
(bandpass sampling, oversampling, reconstruction), the mathematics of quantization and
companding, the clever one-bit converters that live in every phone and audio chip, and the
performance limits of real data converters. It then tells the story of the first great
digital network---the pulse-code-modulated telephone system of 64\,kb/s channels, T1 and E1
frames and multiplexing hierarchies---whose design choices still shape the Internet's
backbone.
\end{chapterintro}

\section{From waveforms to numbers}

In 1937 Alec Reeves, an engineer at International Telephone and Telegraph's Paris laboratory,
filed a patent on a strange idea: instead of transmitting a voice waveform, transmit a
sequence of numbers describing it. Each number would be sent as a short train of on--off
pulses, and because a receiver only had to decide whether a pulse was present, noise could
be removed completely at every repeater. The idea was called \term{pulse-code modulation}
(PCM). It needed thousands of vacuum tubes per channel and was far ahead of its time; the
first commercial PCM system, Bell's T1 carrier, did not enter service until 1962, once
transistors made it cheap. Today PCM, in one form or another, is how every sound, image and
radio waveform enters a computer.

PCM has three steps: \emph{sampling} (discretising time), \emph{quantization} (discretising
amplitude) and \emph{coding} (writing each quantized value as bits). Sampling, done correctly,
loses nothing. Quantization always loses something, and the art is to lose only what does
not matter.

\section{The sampling theorem}

\subsection{Statement and proof sketch}

\begin{keyidea}[title=The sampling theorem]
A signal $x(t)$ with no energy at frequencies $|f|\ge W$ is completely determined by its
samples $x(nT_s)$ taken at rate $f_s=1/T_s>2W$, and can be reconstructed exactly by
\begin{equation}
x(t)=\sum_{n=-\infty}^{\infty}x(nT_s)\,\sinc\big(f_s(t-nT_s)\big).
\label{eq:ch05:interp}
\end{equation}
\end{keyidea}

The rate $2W$ is the \term{Nyquist rate}. The theorem has many fathers: E.~T.~Whittaker
published the interpolation series in 1915, Harry Nyquist identified $2W$ as the maximum
pulse rate of a telegraph channel in 1928, Vladimir Kotelnikov proved the theorem for
communications in 1933, and Claude Shannon stated it in its modern form in 1949.

The proof is a picture. Multiplying $x(t)$ by an impulse train of period $T_s$ in time
convolves its spectrum with an impulse train of spacing $f_s$ in frequency:
\begin{equation}
X_s(f)=f_s\sum_{k=-\infty}^{\infty}X(f-kf_s).
\label{eq:ch05:replicas}
\end{equation}
The sampled spectrum is the original spectrum plus copies---\emph{images} or
\emph{replicas}---centred on every multiple of $f_s$ (Figure~\ref{fig:ch05_sampling_spectra}).
If $f_s>2W$ the replicas do not overlap, and an ideal low-pass filter with cutoff $f_s/2$
recovers $X(f)$ exactly; its impulse response is the sinc of~\eqref{eq:ch05:interp}. If
$f_s<2W$ the replicas overlap and the sum cannot be separated: information has been
destroyed, not merely hidden.

\mdcfig{ch05_sampling_spectra}{Sampling in the frequency domain. Sampling replicates the
spectrum at multiples of $f_s$. Above the Nyquist rate the replicas separate and the
original can be recovered by low-pass filtering (green band); below it they overlap and the
sum is permanently corrupted.}

\subsection{Aliasing}

The overlap has a name, \term{aliasing}, because a frequency component masquerades as a
different one. A sinusoid at $f_0$ sampled at $f_s$ produces exactly the same samples as one
at $f_0+kf_s$ for any integer $k$, and (for real signals) as one at $kf_s-f_0$. In
Figure~\ref{fig:ch05_aliasing_time} a 9\,Hz cosine sampled at 10 samples per second is
indistinguishable from a 1\,Hz cosine. The wagon wheels that turn backwards in Western films,
the strobe light that freezes a spinning fan, and the moir\'e patterns on a striped shirt on
television are all aliasing.

\mdcfig{ch05_aliasing_time}{Aliasing in time: a 9\,Hz cosine (blue) and a 1\,Hz cosine
(dashed) pass through exactly the same samples at 10 samples per second.}

Once aliasing has happened, no processing can undo it. A converter must therefore be
preceded by an \term{anti-aliasing filter} that removes everything above $f_s/2$ before
sampling. No real filter has a brick-wall response, so practical systems sample faster than
the Nyquist rate and leave a transition band. Telephone speech, band-limited to 300--3400\,Hz,
is sampled at 8\,kHz rather than 6.8\,kHz; CD audio, band-limited to 20\,kHz, is sampled at
44.1\,kHz. The ratio of the sample rate to the Nyquist rate, the \term{oversampling ratio},
is the designer's lever: more oversampling means an easier analog filter and more work for
the digital hardware that follows.

\subsection{Sampling real and complex signals}

A real signal with bandwidth $W$ needs $f_s>2W$ real samples per second. A complex baseband
signal occupying $-W/2<f<W/2$ needs $f_s>W$ complex samples per second: its two real
components, I and Q, each need $f_s>W$, so the total real sample rate is the same $2W$. There
is no free lunch, but there is a practical advantage: with complex sampling the spectrum
need not be symmetric about zero, so a receiver can see a signal and its mirror image as
different things (Chapter~\ref{ch:sdr}).

\subsection{Bandpass sampling}

The theorem as stated is for low-pass signals, but a signal that occupies only a band from
$f_L$ to $f_H=f_L+B$ carries information at a rate governed by $B$, not $f_H$. If $f_s$ is
chosen so that the replicas of the band interleave without overlapping, the band aliases
cleanly down to a lower frequency. The condition is
\begin{equation}
\frac{2f_H}{k}\le f_s\le\frac{2f_L}{k-1},\qquad k=1,2,\dots,\lfloor f_H/B\rfloor.
\end{equation}
This \term{bandpass sampling} (or \emph{undersampling}) lets an ADC digitise an IF directly:
a 70\,MHz IF with 20\,MHz bandwidth can be sampled at about 61\,MS/s. The ADC's track-and-hold
must still have analog bandwidth above $f_H$, its aperture jitter must suit the high input
frequency (Section~\ref{sec:ch05:adc}), and the anti-aliasing filter becomes a bandpass
filter that must reject all the other bands that would alias onto the same place.

\subsection{Reconstruction and the zero-order hold}

The sinc interpolator of~\eqref{eq:ch05:interp} is not causal. Real digital-to-analog
converters hold each sample constant for one period, a \term{zero-order hold} (ZOH), and
follow it with an analog \term{reconstruction filter} that removes the images. The hold is a
convolution with a rectangle of width $T_s$, so the output spectrum is multiplied by
$\sinc(f/f_s)$: a gentle roll-off of 3.9\,dB at $f_s/2$ (Figure~\ref{fig:ch05_reconstruction}).
Designs compensate with a digital pre-emphasis filter, oversample so that the band of
interest sits where the droop is small, or both. Every transmitter in this book---including
the AD9364 in the B200---has a ZOH DAC and interpolating digital filters in front of it.

\mdcfig{ch05_reconstruction}{Left: ideal sinc interpolation passes smoothly through the
samples, while a zero-order hold produces a staircase. Right: the staircase's frequency
response $\sinc(f/f_s)$ droops by 3.9\,dB at the Nyquist frequency.}

\section{Quantization}

\subsection{The uniform quantizer}

A $b$-bit quantizer maps each input sample to one of $L=2^b$ levels. The uniform quantizer
spaces the levels by a step $\Delta=2X_{\max}/L$ over a full-scale range $\pm X_{\max}$. The
difference between the output and the input is the \term{quantization error}
$e=Q(x)-x$, bounded by $\pm\Delta/2$ as long as the input does not exceed full scale
(Figure~\ref{fig:ch05_quantizer}). Outside that range the quantizer \term{clips} and the error
grows without bound.

\mdcfig{ch05_quantizer}{A 3-bit mid-rise quantizer, its error sawtooth, and the SQNR of a
full-scale sine versus word length. The simulated values match the $6.02b+1.76$\,dB rule to a
fraction of a decibel.}

When the input is busy---it spans many levels and is not synchronised to the sampling---the
error behaves like a random variable uniformly distributed on $[-\Delta/2,\Delta/2]$,
uncorrelated with the signal and white across $[0,f_s/2]$. Its power is
\begin{equation}
\sigma_e^2=\frac{\Delta^2}{12}.
\end{equation}
For a full-scale sine of amplitude $X_{\max}=2^{b-1}\Delta$ the \term{signal-to-quantization-noise
ratio} is then
\begin{equation}
\text{SQNR}=\frac{X_{\max}^2/2}{\Delta^2/12}=\frac32\,2^{2b}
\quad\Longrightarrow\quad \text{SQNR}_{\text{dB}}=6.02\,b+1.76.
\label{eq:ch05:sqnr}
\end{equation}
This is the most frequently quoted formula in data conversion: \emph{each bit is worth 6\,dB}.
A 12-bit converter, such as the B200's, offers 74\,dB; a 16-bit audio converter 98\,dB; a
24-bit audio converter a theoretical 146\,dB, far below its analog noise floor.

Two refinements matter in practice. First, real signals do not sit at full scale. Speech,
music and OFDM all have large peak-to-average ratios, so the designer must \emph{back off} the
average level to avoid clipping, and every decibel of back-off costs a decibel of SQNR. For a
Gaussian-like signal clipped at four standard deviations, the SQNR is about $6.02b-7.3$\,dB.
Second, if the signal is narrowband compared with $f_s/2$, digital filtering after the ADC
removes the out-of-band quantization noise, giving a \term{processing gain} of
$10\log_{10}(f_s/2B)$\,dB. An SDR that samples at 56\,MS/s and filters to 200\,kHz gains
21.5\,dB, which is why a 12-bit radio can receive a signal 90\,dB below full scale.

\subsection{Dither}

When the input is small or slowly varying, the error is no longer noise-like; it is a
deterministic function of the signal, heard as harmonic distortion or seen as contour bands
in images. Adding a small random \term{dither} signal before quantization---typically noise of
about one LSB---decorrelates the error from the signal. Paradoxically, adding noise improves
quality: the distortion becomes a benign noise floor, and sub-LSB signals remain detectable
in the averaged output. High-resolution measurement ADCs add dither internally.

\section{Non-uniform quantization and companding}

\subsection{The problem with speech}

Speech has a dynamic range of about 40\,dB between loud and soft talkers, and its amplitude
distribution is sharply peaked near zero (approximately Laplacian). A uniform quantizer
wastes its fine levels on large amplitudes that rarely occur and gives quiet talkers a poor
SQNR. The ideal is a quantizer whose step grows in proportion to the signal
amplitude, giving a roughly constant SQNR regardless of level. A logarithmic
characteristic does exactly that.

\subsection{The \texorpdfstring{$\mu$}{mu}-law and A-law}

\term{Companding} (compressing then expanding) implements the logarithmic quantizer as a
nonlinear compressor followed by a uniform quantizer, with the inverse expander at the
receiver. North America and Japan standardised the \term{$\mu$-law}
\begin{equation}
F(x)=\operatorname{sgn}(x)\,\frac{\ln(1+\mu|x|)}{\ln(1+\mu)},\qquad |x|\le1,\ \mu=255,
\end{equation}
and the rest of the world the \term{A-law} with $A=87.6$, which is linear near zero and
logarithmic above $1/A$. Both are specified in ITU-T Recommendation G.711 (1972) as
piecewise-linear approximations with eight segments per polarity, which made them easy to
implement in the digital logic of the time.

Figure~\ref{fig:ch05_companding} shows the result. With a speech-like input, 8-bit $\mu$-law
holds an SQNR of 35 to 38\,dB over a 40\,dB range of input levels, while an 8-bit uniform
quantizer degrades 1\,dB per decibel of level. Matching $\mu$-law across that range with a
uniform quantizer would need about 12 to 13 bits. Eight bits at 8000 samples per second gives
the fundamental rate of digital telephony: \textbf{64\,kb/s}, the \term{DS0} channel.

\mdcfig{ch05_companding}{Left: the $\mu$-law compressor. Right: SQNR against input level for a
speech-like (Laplacian) signal. Eight-bit $\mu$-law keeps the SQNR nearly constant over a 40\,dB
range, performing like a 12-bit uniform quantizer for quiet talkers.}

\subsection{Optimal quantizers}

For a known input distribution, the quantizer that minimises mean-square error satisfies two
conditions found by Stuart Lloyd (1957, published 1982) and Joel Max (1960): each decision
threshold lies midway between adjacent output levels, and each output level is the centroid
(conditional mean) of its decision interval. Iterating the two conditions is the
\term{Lloyd--Max algorithm}. Its multidimensional generalisation, the $k$-means algorithm of
machine learning, designs vector quantizers for speech and image codecs
(Chapter~\ref{ch:sourcecoding}). An important asymptotic result: at high resolution, the
entropy of a uniform quantizer's output is within 0.25\,bit per sample of the best possible
rate for the same distortion, so uniform quantization followed by entropy coding is nearly
optimal---the design behind almost every modern codec.

\section{Differential and predictive coding}

Adjacent speech samples at 8\,kHz are strongly correlated. Rather than quantize each
sample, \term{differential PCM} (DPCM) predicts the next sample from past reconstructed samples
and quantizes only the prediction error, which has much smaller variance. The gain equals the
ratio of signal variance to prediction-error variance, typically 6 to 10\,dB for speech with
a short linear predictor---worth one to two bits per sample. \term{Adaptive DPCM} (ADPCM) also
adapts the quantizer step and predictor to the local statistics. ITU-T G.726 ADPCM at
32\,kb/s delivers quality close to 64\,kb/s PCM and was used to double the capacity of
international circuits and in DECT cordless phones.

Note the structure: the predictor at the encoder uses \emph{reconstructed} samples, exactly what
the decoder will have, so that quantization errors do not accumulate. This ``closed-loop''
prediction is the same principle that makes video codecs and decision-feedback equalizers
(Chapter~\ref{ch:equalization}) stable.

\subsection{Delta modulation}

Take DPCM to the extreme: sample much faster than Nyquist, predict each sample as the
previous reconstructed value, and quantize the error with a single bit. A \term{delta
modulator} transmits 1 if the input is above the staircase approximation and 0 if below,
and the staircase steps by $\pm\delta$ accordingly (Figure~\ref{fig:ch05_delta_mod}). It is
trivially simple and robust to bit errors, but it suffers two impairments. \emph{Granular
noise} appears when the input is flat and the staircase hunts around it. \emph{Slope
overload} appears when the input changes faster than $\delta f_s$ per second and the staircase
cannot keep up. Adaptive delta modulation (CVSD, continuously variable slope delta), which
grows the step during runs of identical bits, found long service in military voice and in
early Bluetooth headsets.

\mdcfig{ch05_delta_mod}{Delta modulation: the staircase follows a slow input closely (with
granular noise) but cannot follow the steep segment on the right (slope overload).}

\section{Oversampling and sigma-delta conversion}

Oversampling spreads the fixed quantization noise power $\Delta^2/12$ over a wider band,
so a digital filter that keeps only the signal band gains 3\,dB per doubling of the
oversampling ratio (OSR)---half a bit per octave. That is too slow to be useful by itself. The
\term{sigma-delta} ($\Sigma\Delta$) modulator changes the game by placing the quantizer inside
a feedback loop with an integrator. The loop forces the \emph{average} of the coarse
quantizer's output to track the input, and it filters the quantization error through a
high-pass \term{noise transfer function}:
\begin{equation}
Y(z)=z^{-1}X(z)+(1-z^{-1})^L E(z)
\end{equation}
for an $L$th-order modulator. The noise is pushed out of the signal band to high frequencies
(Figure~\ref{fig:ch05_sigma_delta}), where a digital decimation filter removes it. The in-band
SQNR now improves by $(6L+3)$\,dB per octave of OSR: a second-order modulator with an OSR of
64 and a one-bit quantizer achieves about 85\,dB. Because a one-bit quantizer is perfectly
linear (two points always define a line), $\Sigma\Delta$ converters avoid the component matching
that limits other architectures. They are in every audio codec, most precision measurement
ADCs, and---as continuous-time bandpass $\Sigma\Delta$ converters---in many cellular receivers.

\mdcfig{ch05_sigma_delta}{Spectra of first- and second-order one-bit sigma-delta modulators
with a tone input. Quantization noise is shaped away from low frequencies; the second-order
loop leaves far less noise in the signal band.}

\section{Real data converters}
\label{sec:ch05:adc}

\subsection{Architectures}

\begin{center}
\begin{tabular}{llll}
\toprule
Architecture & Typical resolution & Typical rate & Where it is used \\
\midrule
Flash & 4--8 bits & up to tens of GS/s & Optical links, oscilloscopes, sub-ranging cores \\
Pipeline & 10--16 bits & 50\,MS/s--1\,GS/s & Base stations, SDRs, instrumentation \\
SAR & 8--18 bits & 1\,kS/s--2\,GS/s & Low power sensors, time-interleaved RF-ADCs \\
Sigma-delta & 16--32 bits & audio to $\approx$100\,MHz bandwidth & Audio, precision, cellular RX \\
Time-interleaved & 8--14 bits & up to 100+\,GS/s & Direct RF sampling, 5G, radar, 800G optics \\
\bottomrule
\end{tabular}
\end{center}

A \emph{flash} converter compares the input with $2^b-1$ thresholds at once; it is the fastest
architecture but its size doubles with every bit. A \emph{successive-approximation} (SAR)
converter performs a binary search with one comparator and a capacitor DAC, and has become
the most power-efficient design at moderate rates. \emph{Pipeline} converters resolve a few
bits per stage and pass the amplified residue on. \emph{Time interleaving} runs $M$ slower
converters in rotation, which multiplies the rate but introduces spurs from mismatched gains,
offsets and timing that must be calibrated digitally. Modern RF-sampling converters, such as
those in AMD's RFSoC devices, sample directly at several gigasamples per second and perform the
mixing and decimation in on-chip digital logic, eliminating the analog mixer altogether.

\subsection{Specifications that matter}

The \term{effective number of bits} (ENOB) inverts equation~\eqref{eq:ch05:sqnr}: measure the
signal-to-noise-and-distortion ratio (SINAD) with a full-scale sine and compute
$\text{ENOB}=(\text{SINAD}_{\text{dB}}-1.76)/6.02$. A 14-bit RF ADC might deliver 11 effective
bits. The \term{spurious-free dynamic range} (SFDR) is the ratio of the signal to the largest
spur, and in a receiver with a strong blocker it often matters more than noise, because a spur
can land on a weak wanted signal.

At high input frequencies, \term{aperture jitter}---random variation of the sampling
instant---dominates. A sinusoid of frequency $f$ has maximum slope $2\pi fA$, so a timing error
$\sigma_t$ produces an amplitude error that limits the SNR to
\begin{equation}
\text{SNR}_{\text{jitter}}=-20\log_{10}(2\pi f\sigma_t)\ \text{dB}.
\end{equation}
With 100\,fs of jitter, the limit at 1\,GHz is 64\,dB, barely 10 bits, whatever the converter's
resolution (Figure~\ref{fig:ch05_jitter}). This is why clock generation is a serious design
discipline in every RF-sampling system, and why the B200 and every base station specify the
phase noise of their reference oscillators.

\mdcfig[0.7]{ch05_jitter}{The SNR ceiling set by sampling-clock jitter. At gigahertz input
frequencies even 50\,fs of jitter limits a converter to about 10 to 11 effective bits.}

\section{Digital telephony}

\subsection{The 64 kb/s channel}

By the late 1950s, the Bell System faced a crisis in its local networks: the cable ducts
between exchanges in big cities were full. Replacing copper pairs was expensive; making each
pair carry more calls was not. The answer was T1 carrier, which put 24 PCM voice channels on
two twisted pairs, with regenerative repeaters every 6000 feet (the spacing of the existing
loading coils, whose manholes were already there). Each repeater received a noisy stream of
pulses, decided what each bit was, and retransmitted a clean pulse. For the first time in
history, a long-distance voice circuit had the same quality as a short one.

\subsection{Time-division multiplexing and the T1 and E1 frames}

The channels share the line by \term{time-division multiplexing} (TDM): each channel gets
a time slot of eight bits in a repeating frame. The T1 (DS1) frame (Figure~\ref{fig:ch05_t1_frame})
contains 24 slots plus one framing bit, 193 bits every 125\,$\mu$s, giving
$193\times8000=1.544$\,Mb/s. The framing bits follow a fixed pattern over 12 frames (the
superframe) or 24 frames (the extended superframe, which also carries a CRC-6 and a
4\,kb/s data link). Signalling (on-hook, off-hook, dial pulses) was originally carried by
stealing the least significant bit of each channel in every sixth frame---``robbed-bit''
signalling---which is why modem data over digital trunks was limited to 56\,kb/s instead of
64. The European E1 frame, standardised by CEPT, uses 32 slots of eight bits: slot 0 for
framing and alarms, slot 16 for signalling, and 30 for voice, for 2.048\,Mb/s. Its cleaner
design carries 64\,kb/s ``clear channels'' without robbed bits.

\mdcfig{ch05_t1_frame}{The T1 (DS1) frame: one framing bit and 24 eight-bit $\mu$-law samples
every 125\,$\mu$s, for 1.544\,Mb/s.}

\subsection{The multiplexing hierarchy}

Frames are multiplexed again into higher rates. The North American \term{plesiochronous digital
hierarchy} (PDH) combined four DS1s into a DS2 (6.312\,Mb/s) and seven DS2s into a DS3
(44.736\,Mb/s, 672 voice channels); the European hierarchy went 2.048, 8.448, 34.368 and
139.264\,Mb/s. ``Plesiochronous'' means \emph{nearly} synchronous: each tributary has its own
clock, so the multiplexer inserts \emph{justification} (stuffing) bits to absorb small rate
differences. Extracting one DS1 from a DS3 required demultiplexing everything.

The \term{synchronous optical network} (SONET, 1988) and synchronous digital hierarchy (SDH)
replaced this with byte-synchronous frames of 125\,$\mu$s at 51.84\,Mb/s (STS-1) and multiples
(OC-3 at 155.52\,Mb/s, OC-48 at 2.488\,Gb/s, OC-192 at 9.953\,Gb/s), with pointers that locate
each tributary directly. SONET's 125\,$\mu$s frame---the sampling period of an 8\,kHz voice
channel---carried the Internet's backbone traffic through the 1990s and 2000s, and the timing
distribution networks built for it still synchronise cellular base stations today. The story
continues in Chapter~\ref{ch:wireline}.

\subsection{Line coding on the wire}

A copper pair cannot carry DC through its transformers, and the receiver must recover the
bit clock from the data. T1 therefore uses \term{alternate mark inversion} (AMI): zeros are
sent as no pulse and ones as pulses of alternating polarity, which gives a spectrum with no DC
and a built-in error check (two consecutive pulses of the same polarity are a
\emph{bipolar violation}). Long runs of zeros would starve the clock recovery, so T1 requires
a minimum ones density, later replaced by \emph{B8ZS} substitution, which replaces eight zeros
with a pattern containing deliberate violations that the receiver recognises. E1 uses the
similar HDB3. Line codes are treated fully in Chapter~\ref{ch:baseband}.

\subsection{Echo, delay and the end of the circuit}

Digital telephony exposed problems analog networks had hidden. The \term{hybrid} that converts
the two-wire subscriber loop to the four-wire transmission network leaks some of the
far-end's voice back as echo. On short circuits the echo arrives so quickly that it merges
with sidetone, but on satellite circuits (500\,ms round trip, Chapter~\ref{ch:satellite}) or
packet networks it is intolerable. \term{Adaptive echo cancellers}, LMS filters that learn the
echo path and subtract a replica (Chapter~\ref{ch:equalization}), were among the first mass
applications of adaptive signal processing and are in every telephone network and speakerphone
today.

From the 2000s, voice moved from circuits to packets. Voice over IP (VoIP) and voice over LTE
(VoLTE) still sample at 8 or 16\,kHz, but compress the samples with codecs such as AMR-WB, EVS
and Opus at 6 to 64\,kb/s (Chapter~\ref{ch:sourcecoding}). The 125\,$\mu$s frame survives only
in legacy interfaces; the 64\,kb/s G.711 codec survives as the universal fallback that every
VoIP system still supports.

\begin{history}[title=The 8 kHz heartbeat]
The number 8000 appears everywhere in communications because telephone speech was
band-limited to 3.4\,kHz and the designers of T1 chose a comfortable margin above the
6.8\,kHz Nyquist rate. From that choice follow the 125\,$\mu$s frame, the 64\,kb/s DS0, the
1.544 and 2.048\,Mb/s primary rates, SONET's 51.84\,Mb/s, the 8\,kHz frame clock of
ISDN and GSM's 20\,ms speech frames. A sampling-rate decision made for a vacuum-tube
filter in the 1950s still sets the timing of the global network.
\end{history}

\begin{pitfall}[title=Clipping is worse than noise]
Quantization noise degrades gracefully; clipping does not. A signal that clips even
occasionally generates broadband distortion and, in a transmitter, spectral regrowth into
adjacent channels. For OFDM signals with 10\,dB or more of peak-to-average ratio, the rule of
thumb is to keep the RMS level 10 to 12\,dB below full scale in the ADC and DAC, and to budget
the lost SQNR accordingly. In the B200, setting the digital TX amplitude above about 0.7 is a
common way to generate a very clean spectral splatter.
\end{pitfall}

\begin{labbox}[title=Lab: sampling and quantization]
The notebook \lab{labs/ch05\_sampling\_quantization.ipynb} lets you sweep an input frequency
through the Nyquist zones and watch it alias; compares sinc, linear and zero-order-hold
reconstruction; measures SQNR against bits, back-off and dither; implements $\mu$-law and A-law
companding and reproduces Figure~\ref{fig:ch05_companding}; and builds first- and second-order
sigma-delta modulators whose output spectra update as you change the OSR.
\end{labbox}

\problems
\begin{enumerate}[label=\thechapter.\arabic*]
\item A signal contains components at 3, 7 and 12\,kHz and is sampled at 16\,kHz without an
anti-aliasing filter. At what frequencies do the components appear after ideal reconstruction?
\item Find all valid bandpass sampling rates for a signal occupying 70--80\,MHz. Which is the
lowest? What practical margin would you add?
\item Derive the ZOH frequency response $\sinc(f/f_s)$ and design a three-tap FIR
compensation filter that flattens it to within 0.2\,dB up to $0.4f_s$.
\item A 14-bit ADC samples at 125\,MS/s. What is the ideal SQNR? If the signal of interest is a
1\,MHz channel after digital filtering, what is the in-band SQNR? What jitter would limit it at
an input frequency of 200\,MHz?
\item Show that for a uniform quantizer with a Gaussian input of standard deviation $\sigma$
clipped at $\pm4\sigma$, SQNR $\approx 6.02b-7.3$\,dB (neglect the clipping noise).
\item Implement the Lloyd--Max algorithm for a 3-bit quantizer and a Laplacian input. Compare the
MSE with that of the best uniform quantizer.
\item Compute the bit rates of a T1 superframe's framing pattern, a DS3 and an OC-12. How many
64\,kb/s channels does an OC-192 carry?
\item Explain why a first-order sigma-delta modulator with a DC input produces periodic
patterns (idle tones), and how dither or higher order removes them.
\item \emph{(Simulation)} Record your voice at 48\,kHz (or use any speech file), decimate to
8\,kHz, and encode it with 8-bit uniform PCM, 8-bit $\mu$-law and 4-bit ADPCM. Compare SQNR
and listen to the results.
\end{enumerate}

\furtherreading
\begin{itemize}
\item N.~S.~Jayant and P.~Noll, \emph{Digital Coding of Waveforms}, Prentice Hall, 1984. The
classic on PCM, DPCM, ADPCM and delta modulation.
\item R.~Schreier and G.~C.~Temes, \emph{Understanding Delta-Sigma Data Converters}, 2nd ed.,
Wiley-IEEE, 2017.
\item W.~Kester (ed.), \emph{Data Conversion Handbook}, Analog Devices / Newnes, 2005 (free online).
\item B.~M.~Oliver, J.~R.~Pierce and C.~E.~Shannon, ``The philosophy of PCM,''
\emph{Proc. IRE}, 36, 1948. A remarkably modern paper.
\item J.~Bellamy, \emph{Digital Telephony}, 3rd ed., Wiley, 2000. T1, E1, SONET and the
engineering of the digital network.
\item ITU-T Recommendations G.711 (PCM), G.726 (ADPCM), G.704 (frame structures).
\end{itemize}
